\chapter{Background}

\section{The pipeline}
A benchmark circuit from MQT Bench is transpiled to the gate set $\{R_z,\sqrt X,X,\mathrm{CX}\}$ and cut into layers by an
as-soon-as-possible schedule of its dependency graph. A qubit mapper (HQA, Hungarian-assignment based) places the logical
qubits on cores layer by layer. Whenever the placement changes between two layers, the qubit has to move, and the move is expanded into one
\emph{hop} per link along a path found by Dijkstra's algorithm on the core graph: shortest in hops, with the product of
link reliabilities as the tie-break. The result is a \emph{trace}: for every layer, a list of hops
$(\text{qubit},\text{source core},\text{destination core})$. Every hop of the traces used here joins two adjacent cores.
The traces carry no gate-level dependencies, so a layer boundary stands in for them.

ACE then executes the trace: each hop is a request for one EPR pair between two adjacent cores, submitted through
SeQUeNCe's resource-reservation protocol (RSVP), which books one communication memory at each end for a time window.

\section{ACE and the native SeQUeNCe backend}
ACE adds, on top of SeQUeNCe, adaptive continuous pre-generation of entanglement (cached pairs in spare memories), a compiler
controller that pre-generates exactly the pairs the trace will need (the policies \emph{fixed} and \emph{dynamic}),
and a request manager that submits a layer's requests together. The native backend instead uses a central scheduler
with one generation job per link. The two backends are never compared in this report.

\section{Policies}
\emph{On-demand} generates a pair only when a request asks for it, so it isolates contention. \emph{Dynamic} compiler
pre-generation prepares pairs ahead of the layer that needs them, using at most \code{budget - 1} memories per core as a
compiler bank and leaving one for demand. Autonomous caches (\emph{continuous}, \emph{adaptive-continuous}) guess which pairs
will be needed.

\section{Parameters in force}
Unless stated, the physical parameters are ACE's configuration: memory fidelity $0.95$, memory efficiency $0.95$, coherence
time $1$\,s, detector efficiency $0.95$, gate and measurement fidelity $0.99$, classical delay $100\,\mu$s plus $20\,\mu$s
per hop, and a router-to-BSM fibre of $0.025$ (read as metres) with attenuation $0.2$\,dB/km. The classical delays, detector efficiency, fidelities and attenuation match the metro-scale adaptive-continuous protocol paper; the link length, the coherence time and the memory efficiency differ from it, and the last two also from the group's multi-core paper (Appendix~\ref{app:params}).
